\section{Related Work}
\label{sec:relatedwork}

Optimizing LLM inference is an active area of research ~\cite{helix, griggs2024melangecostefficientlarge, miao2023spotserveservinggenerativelarge, zhu2024nanoflowoptimallargelanguage, kossmann2024gpuhalfemptyhalffullpractical, strati2024dejavukvcachestreamingfast, 
lee2024infinigenefficientgenerativeinference, singhania2024lokilowrankkeysefficient, zheng2024sglangefficientexecutionstructured, ma2024compressingkvcachelongcontext,mahmud2024mapleframeworkactivepreference}. Various techniques have been proposed to improve different aspects of LLM serving like batching~\cite{orca, vllmsosp, sarathiserve2024}, disaggregation~\cite{patel2023splitwise, distserve2024, tetriinfer}, scheduling~\cite{flexgen, fastserve}. However, central to all these techniques is the need for efficient \kvcache memory management. Since \vllm, \pa has been adopted in various serving frameworks e.g., TensorRT-LLM~\cite{trtllmgithub}, LightLLM~\cite{lightllm:github}, and libraries e.g., \flashattention~\cite{fagithub} and \flashinfer~\cite{figithub}. Some other concurrent works have also proposed optimizations for attention kernels~\cite{flashattnhopper,mirage-tensor-compiler, thunderkittens, h2o2023, flashdecoding++,leanattention,flash-attention-3, pod-attn-2024}. Deploying new kernels for inference is hard with \pa.  \sysname offers an alternate -- which we believe is also a more principled -- approach to dynamic \kvcache memory management that makes it easier to deploy new kernels.

In a recent work, GMLake~\cite{gmlakeasplos24} showed that using CUDA virtual memory support can mitigate fragmentation in DNN training jobs, increasing training batch size. In particular, GMLake uses CUDA support to coalesce multiple smaller physical memory pages into a single virtually contiguous object that can prevent out-of-memory errors. In contrast, \sysname targets optimizing inference workloads. %Different from training, LLM inference is latency sensitive and requires smaller granularity allocations. We proposed various LLM inference specific optimizations to meet these requirements. 

%Similar to our approach, PyTorch also recently added support for expandable segments~\cite{expandable-segments} that dynamically attach physical memory pages to a pre-reserved virtual memory buffer. PyTorch uses expandable segments opportunistically, using synchronous memory allocation and only in multiples of 2MB granularity. In contrast, \sysname caters to the specific requirements of the \kvcache.  Some other concurrent works have also proposed optimizations for attention kernels~\cite{flashattnhopper,mirage-tensor-compiler, thunderkittens, h2o2023, flashdecoding++,leanattention,flash-attention-3, pod-attn-2024}. As one would expect, most of these kernels are not compatible with \pa. \sysname would make it easier to deploy them.